% Generated draft: requires mathematical and source review.
\documentclass[11pt]{article}
\usepackage[margin=25mm]{geometry}
\usepackage{fontspec}
\setmainfont{Latin Modern Roman}
\usepackage{amsmath,amssymb,mathtools,unicode-math}
\setmathfont{Latin Modern Math}
\usepackage{enumitem,needspace,xurl,hyperref}
\setlength{\parindent}{0pt}
\setlength{\parskip}{5pt}
\setlength{\emergencystretch}{3em}
\newcommand{\N}{\mathbb{N}}
\newcommand{\Z}{\mathbb{Z}}
\newcommand{\Q}{\mathbb{Q}}
\newcommand{\R}{\mathbb{R}}
\newcommand{\C}{\mathbb{C}}
\newcommand{\F}{\mathbb{F}}
\newcommand{\A}{\mathbb{A}}
\newcommand{\PP}{\mathbb{P}}
\newcommand{\E}{\mathbb{E}}
\newcommand{\eps}{\varepsilon}
\newcommand{\dd}{\,\mathrm{d}}
\begin{document}
\section*{UnsolvedMath ID: 30000878}
\textbf{Source problem code:} catalogue entry
\subsection{Definitions and mathematical statement}
This entry is marked unrecoverable because it is the workshop's broad topic overview rather than a discrete mathematical problem. Its clean statement is null; detailed questions belong to subsequent records, whose identifiers are not supplied here.

For a group $G$, the integral group ring $\Z[G]$ consists of finite sums $\sum_g a_gg$ with convolution multiplication, and $U(\Z[G])$ denotes its invertible elements. A torsion unit has finite multiplicative order; a torsion subgroup is a subgroup all of whose elements have finite order. The retained overview lists:
\begin{enumerate}
\item constructing finite-index subgroups of the unit group and studying subgroups generated by units of prescribed types;
\item determining properties of unit groups, especially for infinite groups;
\item describing torsion units and torsion subgroups for finite nonsoluble groups and for infinite groups;
\item resolving the first Zassenhaus conjecture;
\item determining whether a finite $p$-group is determined up to group isomorphism by its group algebra over $\F_p$.
\end{enumerate}
The last topic can be recognized as a group-algebra reconstruction question, but no independent theorem is promoted from this overview. The unit types, properties, group classes, normalization of units, and exact Zassenhaus formulation are not specified and should be taken from the detailed source questions, rather than supplied by name alone.

The dated assessment reports a counterexample to the first Zassenhaus conjecture and substantial special cases of other topics, while describing the collection as unresolved. This is retained as supplied literature context, not a claim that every historically listed question remains open. In particular, the broad partially-solved label cannot assign a reliable present-day status to each topic individually.
\subsection{Short English statement}
This provenance-only entry lists several integral-group-ring research themes, including unit groups, torsion units, Zassenhaus, and modular group-algebra reconstruction. It is not a single conjecture, and its bundled status does not determine the current status of each listed question.
\subsection{Sources}
\begin{itemize}
\item Local UnsolvedMath registry, permanent ID 30000878, code catalogue entry, ``Structural Problems for Units of Integral Group Rings.'' Proposed year 2007; proposer unspecified.
\item \emph{Mini-Workshop: Arithmetik von Gruppenringen} (2008), Oberwolfach Reports, pp.\ 3209--3240, DOI 10.4171/owr/2007/55; \url{https://doi.org/10.4171/owr/2007/55}. The displayed report year 2008 differs from the DOI and proposed year 2007.
\item Supplied ``counterexample to first Zassenhaus conjecture'': \url{https://doi.org/10.1016/j.aim.2018.10.004}.
\item Supplied ``prime graph reduction'': \url{https://arxiv.org/abs/1612.05025}.
\item Statement status unrecoverable; clean statement null. The supplied official-report verification identifies a workshop-topic overview and says detailed questions appear in subsequent records.
\item Literature triage checked 2026-08-21: partially solved. It reports a Zassenhaus counterexample and other special cases but no complete resolution of the collection. It gives no separate theorem/status accounting for each listed topic. Created 2026-08-25; updated 2026-09-16. Source and current-status claims require independent verification.
\item Registry curation by Ulam AI / UnsolvedMath Contributors, CC BY 4.0; underlying report and literature provenance remain as supplied. Unreviewed provenance-only draft.
\end{itemize}
\end{document}
